\section{Targeted and margin-guided tests at small budgets (Studies 2 and 3)}\label{sec:strategies}

This section is post hoc and exploratory: its tests were designed after the results of Studies~2 and~3 were known.
For the full-domain DAS alignments of Studies~2 and~3, we computed from the censuses and the saved networks the expected
share of alignments that each of six tests detects at 12 to 120 model calls, charging 3 calls per single query, 4 per
joint query and 1 per gradient (\cref{tab:strategies}):
\begin{itemize}
  \item the alternating single test and the joint test of the frozen comparisons;
  \item an adaptive single test, which alternates for the first four queries and then spends the rest on the variable
  with more observed failures, alternating on ties;
  \item the same with margins: after four queries without a failure, it spends the rest on the variable whose two
  queries had the smaller mean margin;
  \item an oracle-targeted single test, which spends every query on whichever variable fails more often, as a tester
  could do only by knowing it in advance;
  \item a margin-guided joint test, which draws joint-panel queries, runs the two single queries that each of them
  reduces to together with the gradients that $\sigma$ needs, scores every joint-panel query that the drawn single
  queries compose, and runs the top-scored ones. A candidate costs 8 calls in Study~2 (two single queries and two
  gradients) and 7 in Study~3 (one gradient); at 120 calls the test draws \StTwoGuidedCandidates{} candidates and runs
  \StTwoGuidedRuns{} joint queries in Study~2, and \StThreeGuidedCandidates{} and \StThreeGuidedRuns{} in Study~3. We
  estimated it from \StReplicates{} simulated tests per alignment and budget, beside the same test with the
  candidates run in random order.
\end{itemize}

\begin{table}[t]
\centering
\small
\setlength{\tabcolsep}{5pt}
\begin{tabular}{@{}l c c c c c c c c@{}}
\toprule
& \multicolumn{4}{c}{Study~2 (model calls)} & \multicolumn{4}{c}{Study~3 (model calls)} \\
\cmidrule(lr){2-5}\cmidrule(lr){6-9}
Test & 12 & 24 & 60 & 120 & 12 & 24 & 60 & 120 \\
\midrule
Alternating single & 0.629 & 0.723 & 0.788 & 0.816 & 0.627 & 0.765 & 0.874 & 0.911 \\
Adaptive single (margins) & 0.629 & 0.723 & 0.787 & 0.815 & 0.627 & 0.796 & 0.881 & 0.903 \\
Oracle-targeted single & 0.644 & 0.732 & 0.792 & 0.818 & 0.740 & 0.845 & 0.910 & 0.927 \\
Joint & 0.527 & 0.648 & 0.746 & 0.788 & 0.676 & 0.808 & 0.897 & 0.924 \\
Margin-guided joint & 0.481 & 0.623 & 0.757 & 0.799 & 0.489 & 0.653 & 0.839 & 0.901 \\
Margin-guided, random order & 0.481 & 0.623 & 0.754 & 0.796 & 0.488 & 0.652 & 0.834 & 0.897 \\
\bottomrule

\end{tabular}
\caption{Expected share of the full-domain DAS alignments detected (post hoc). The adaptive single test that switches
on observed failures is identical to the alternating test at every budget and is not shown. The margin-guided rows
are means of \StReplicates{} simulated tests per alignment (standard error at most \StMaxSe).}
\label{tab:strategies}
\end{table}

\paragraph{Adapting to observed failures changes nothing.}
The adaptive single test can switch only after it has observed a failure, and a test that has observed a failure has
already detected the alignment; its detection probability is therefore the alternating test's at every budget (the
two agree to \StAdaptiveDiff). Adapting to margins can switch before any failure. In Study~3, where $B$ is the weaker
variable, it raises the share detected from \StThreeAltTwentyFour{} to \StThreeAdaptMarginTwentyFour{} at 24 calls
but lowers it from \StThreeAltOneTwenty{} to \StThreeAdaptMarginOneTwenty{} at 120 calls, because two queries per
variable sometimes point to the wrong one; in Study~2 it changes almost nothing. Knowing the weaker variable in advance would help most:
the oracle-targeted single test is the best single test in both studies, and in Study~3 the joint test comes close to
it (\StThreeJointOneTwenty{} against \StThreeOracleOneTwenty{} at 120 calls), as the post hoc analysis of
\cref{sec:multisite} leads one to expect. These tests only detect failures; for the failure rate
itself, \citet{asiaee2026certified} gives confidence sequences that stay valid when a test stops early or, through
importance weights, aims its queries at likely failures.

\paragraph{Margin guidance does not pay at these budgets.}
The margin-guided joint test detects fewer alignments than the alternating single test at every budget in both
studies (\StTwoGuidedOneTwenty{} against \StTwoAltOneTwenty{} in Study~2 and \StThreeGuidedOneTwenty{} against
\StThreeAltOneTwenty{} in Study~3, at 120 calls), and at most slightly more than the same test with its candidates
in random order (\StTwoGuidedRandomOneTwenty{} and \StThreeGuidedRandomOneTwenty). On the joint-only alignments, which no
single test can detect, the score does pick better candidates: at 120 calls the guided test detects
\StTwoJoGuidedOneTwenty\% of Study~2's \StTwoJoCount{} joint-only DAS alignments and \StThreeJoGuidedOneTwenty\% of
Study~3's \StThreeJoCount, against \StTwoJoGuidedRandomOneTwenty\% and \StThreeJoGuidedRandomOneTwenty\% with the
candidates in random order. But a plain joint test, which runs 30 joint queries instead of a handful, detects \StTwoJoJointOneTwenty\% and
\StThreeJoJointOneTwenty\%. Scoring a candidate costs more than running it (7 or 8 calls against 4), so at budgets of
this size the score cannot save queries. It pays where the single census is taken anyway: then every joint query can
be scored for the price of the single queries' gradients (\cref{sec:theory}).
